problem:  
Let \( M^n \) be a compact, connected Riemannian manifold of nonnegative sectional curvature admitting an effective, isometric torus action  
\[
T^k \curvearrowright M,
\]
where \(k \ge \lfloor n/2 \rfloor - 1\).  Denote by \(M^* = M/T^k\) the quotient equipped with its Alexandrov metric.  
Existing classification results describe the topology of \(M\) when \(k\) is maximal (\(k = \lfloor n/2 \rfloor\)), and in many of those cases the orbit space \(M^*\) is known to be a simplex or disk.  

However, very little is currently understood for *nearly maximal* symmetry rank.  

**Question:**  
For \(n \ge 6\), does there exist a compact nonnegatively curved manifold \(M^n\) with symmetry rank  
\[
k = \lfloor n/2 \rfloor - 1
\]
such that the quotient \(M^* = M/T^k\) is **not homeomorphic to a manifold with boundary**?   
Equivalently, can there exist a nonnegatively curved manifold with symmetry rank one less than maximal whose Alexandrov orbit space still contains genuine stratified singularities (nonmanifold points) in codimension ≥2?

Why is it a "good" problem:  
This version pinpoints a sharp, unexplored threshold between known rigidity and potential geometric flexibility. For maximal symmetry rank, the orbit space structure is completely characterized: the quotient is a manifold, and \(M\) is strongly constrained topologically. The next symmetry rank down—just one dimension less—is the frontier where current classification techniques fail. Determining whether nonmanifold quotients can occur here goes straight to the heart of how curvature and symmetry interact to produce or suppress singularities.  

The problem is tightly connected to the Grove symmetry program, Alexandrov geometry, and curvature rigidity theory. It is simple to state, naturally builds on known deep results, yet the answer is genuinely unknown. Resolving it would either reveal a new rigidity phenomenon or exhibit novel geometric examples with controlled singular orbit structure—an outcome of significant conceptual and structural impact.